\chapter{Erste Schritte}
\label{ch:basics/getting-started}

Den grundlegenden Aufbau eines Mod-Verzeichnisses zu verstehen, ist beim
Erstellen von Mods unerlässlich. In diesem Kapitel erfährst du, wie
Modding in Luanti funktioniert, und erstellst deinen ersten Mod.

\section{Was sind Spiele und Mods?}\label{was-sind-spiele-und-mods}

Die Stärke von Luanti liegt darin, dass sich Spiele einfach entwickeln
lassen, ohne eigene Voxel-Grafik, Voxel-Algorithmen oder ausgefeilten
Netzwerkcode schreiben zu müssen.

In Luanti ist ein Spiel eine Sammlung von Modulen, die zusammenarbeiten,
um den Inhalt und das Verhalten des Spiels bereitzustellen. Ein Modul,
allgemein als Mod bezeichnet, ist eine Sammlung von Skripten und
Ressourcen. Es ist möglich, ein Spiel mit nur einem Mod zu erstellen,
aber das wird selten gemacht, weil sich dann Teile des Spiels nicht mehr
so einfach unabhängig voneinander anpassen oder ersetzen lassen.

Ebenso ist es möglich, Mods außerhalb eines Spiels zu verbreiten. Dann
sind sie ebenfalls \emph{Mods}, aber im traditionelleren Sinn:
\emph{Modifikationen}. Diese Mods verändern oder erweitern die
Eigenschaften eines Spiels.

Sowohl die Mods, die in einem Spiel enthalten sind, als auch Mods von
Dritten nutzen dieselbe API (Programmierschnittstelle).

Dieses Buch behandelt die wichtigsten Teile der Luanti-API und richtet
sich gleichermaßen an Spieleentwickler und Modder.

\section{Wo werden die Mods
gespeichert?}\label{wo-werden-die-mods-gespeichert}

Jeder Mod hat sein eigenes Verzeichnis, in dem sich seine
Lua-Quelltexte, Texturen,\index{textur@Textur} Modelle und Tondateien
befinden. Luanti sucht an verschiedenen Orten nach Mods. Diese Orte
werden allgemein \emph{Mod-Ladepfade} genannt.

Für eine bestimmte Welt werden drei Mod-Speicherorte geprüft, und zwar
in dieser Reihenfolge:

\begin{enumerate}
\def\labelenumi{\arabic{enumi}.}
\tightlist
\item
  Spiel-Mods. Das sind die Mods des Spiels, in dem die Welt läuft.\\
  Z.~B.:
  \texttt{minetest/\allowbreak{}games/\allowbreak{}minetest\_\allowbreak{}game/\allowbreak{}mods/\allowbreak{}}
  oder\\
  \texttt{/\allowbreak{}usr/\allowbreak{}share/\allowbreak{}minetest/\allowbreak{}games/\allowbreak{}minetest/\allowbreak{}}
\item
  Globale Mods. Der Ort, an dem Mods fast immer installiert werden. Im
  Zweifelsfall legst du Mods hier ab.\\
  Z.~B.: \texttt{minetest/\allowbreak{}mods/\allowbreak{}}
\item
  Welt-Mods. Der Ort für Mods, die zu einer bestimmten Welt gehören.\\
  Z.~B.:
  \texttt{minetest/\allowbreak{}worlds/\allowbreak{}world/\allowbreak{}worldmods/\allowbreak{}}
\end{enumerate}

\texttt{minetest} ist das
Benutzerdatenverzeichnis.\index{benutzerdatenverzeichnis@Benutzerdatenverzeichnis}
Seinen Ort findest du, indem du Luanti öffnest und in der Registerkarte
„Über`` auf „Benutzerdatenverzeichnis öffnen`` klickst. Wurde Luanti
über Flathub installiert, passiert bei diesem Klick unter Umständen
nichts. Das Verzeichnis liegt dann normalerweise unter
\texttt{\textasciitilde{}/.\allowbreak{}var/\allowbreak{}app/\allowbreak{}org.\allowbreak{}luanti.\allowbreak{}luanti/.\allowbreak{}minetest/\allowbreak{}}.

Beim Laden prüft Luanti die oben genannten Orte der Reihe nach. Stößt es
auf einen Mod mit demselben Namen wie ein zuvor gefundener, wird der
spätere anstelle des früheren geladen. Du kannst also Spiel-Mods
überschreiben, indem du einen gleichnamigen Mod in das globale
Mod-Verzeichnis legst.

\section{Deinen ersten Mod erstellen}\label{deinen-ersten-mod-erstellen}

\subsection{Mod-Verzeichnis}\label{mod-verzeichnis}

Geh in das globale Mod-Verzeichnis (Über \textgreater{}
Benutzerdatenverzeichnis öffnen \textgreater{} mods) und lege einen
neuen Ordner namens \texttt{mymod} an. \texttt{mymod} ist der Name des
Mods.

Jeder Mod sollte einen eindeutigen \emph{Mod-Namen} haben, eine
technische Kennung (ID), über die auf den Mod verwiesen wird. Mod-Namen
dürfen Buchstaben, Ziffern und Unterstriche enthalten. Ein guter Name
beschreibt, was der Mod tut. Das Verzeichnis mit den Bestandteilen eines
Mods muss genauso heißen wie der Mod. Ob ein Mod-Name noch frei ist,
findest du mit einer Suche auf \url{https://content.luanti.org} heraus.

\begin{Shaded}
\begin{Highlighting}[]
\NormalTok{mymod}
\NormalTok{├── textures}
\NormalTok{│   └── mymod\_node.png}
\NormalTok{├── init.lua}
\NormalTok{└── mod.conf}
\end{Highlighting}
\end{Shaded}

Ein Mod benötigt nur eine
\texttt{init.\allowbreak{}lua}\index{init.lua@\texttt{init.\allowbreak{}lua}}.
Eine
\texttt{mod.\allowbreak{}conf}\index{mod.conf@\texttt{mod.\allowbreak{}conf}}
wird jedoch empfohlen, und je nach Funktionsumfang können weitere
Bestandteile nötig sein.

\subsection{mod.conf}\label{mod.conf}

Lege eine Datei \texttt{mod.\allowbreak{}conf} mit folgendem Inhalt an:

\begin{Shaded}
\begin{Highlighting}[]
\NormalTok{name = mymod}
\NormalTok{description = Fügt foo, bar und bo hinzu.}
\NormalTok{depends = default}
\end{Highlighting}
\end{Shaded}

Diese Datei enthält die Metadaten\index{metadaten@Metadaten} des Mods,
darunter seinen Namen, seine Beschreibung und weitere Angaben.

\subsection{init.lua}\label{init.lua}

Lege eine Datei \texttt{init.\allowbreak{}lua} mit folgendem Inhalt an:

\begin{Shaded}
\begin{Highlighting}[]
\FunctionTok{print}\OperatorTok{(}\StringTok{"Diese Datei wird beim Laden ausgeführt!"}\OperatorTok{)}

\VariableTok{core}\OperatorTok{.}\NormalTok{register\_node}\OperatorTok{(}\StringTok{"mymod:node"}\OperatorTok{,} \OperatorTok{\{}
    \VariableTok{description} \OperatorTok{=} \StringTok{"Das ist ein Node"}\OperatorTok{,}
    \VariableTok{tiles} \OperatorTok{=} \OperatorTok{\{} \StringTok{"mymod\_node.png"} \OperatorTok{\},}
    \VariableTok{groups} \OperatorTok{=} \OperatorTok{\{} \VariableTok{cracky} \OperatorTok{=} \DecValTok{1} \OperatorTok{\},}
\OperatorTok{\})}

\VariableTok{core}\OperatorTok{.}\NormalTok{register\_craft}\OperatorTok{(\{}
    \FunctionTok{type} \OperatorTok{=} \StringTok{"shapeless"}\OperatorTok{,}
    \VariableTok{output} \OperatorTok{=} \StringTok{"mymod:node 3"}\OperatorTok{,}
    \VariableTok{recipe} \OperatorTok{=} \OperatorTok{\{} \StringTok{"default:dirt"}\OperatorTok{,} \StringTok{"default:stone"} \OperatorTok{\},}
\OperatorTok{\})}
\end{Highlighting}
\end{Shaded}

Die \texttt{init.\allowbreak{}lua} ist der Einstiegspunkt eines Mods und
wird ausgeführt, wenn der Mod geladen wird.

\subsection{Zusammenfassung}\label{zusammenfassung}

Dieser Mod heißt „mymod``. Er besteht aus zwei Textdateien:
\texttt{init.\allowbreak{}lua} und \texttt{mod.\allowbreak{}conf}. Das
Skript gibt eine Nachricht aus und registriert dann einen Node und ein
Handwerksrezept -- beides wird später erklärt. Es gibt eine einzige
Abhängigkeit, den default-Mod\footnote{\url{https://content.luanti.org/metapackages/default/}},
der normalerweise in Minetest Game enthalten ist. Außerdem liegt in
\texttt{textures/\allowbreak{}} eine Textur für den Node.

\section{Abhängigkeiten}\label{abhuxe4ngigkeiten}

Eine Abhängigkeit besteht, wenn ein Mod einen anderen Mod benötigt, der
vor ihm geladen werden muss. Ein Mod kann darauf angewiesen sein, dass
Code, Gegenstände oder andere Ressourcen eines anderen Mods verfügbar
sind.

Es gibt zwei Arten von Abhängigkeiten: feste und optionale. Bei beiden
wird der andere Mod zuerst geladen. Ist der Mod, von dem abgehangen
wird, nicht verfügbar, führt eine feste Abhängigkeit dazu, dass der
eigene Mod nicht geladen wird, während eine optionale Abhängigkeit
höchstens dazu führt, dass weniger Funktionen zur Verfügung stehen.

Eine optionale Abhängigkeit ist nützlich, wenn du einen anderen Mod
optional unterstützen möchtest; so lassen sich zusätzliche Inhalte
aktivieren, wenn der Benutzer beide Mods gleichzeitig verwendet.

Abhängigkeiten werden in der \texttt{mod.\allowbreak{}conf} als
kommagetrennte Liste angegeben:

\begin{Shaded}
\begin{Highlighting}[]
\NormalTok{depends = modeins, modzwei}
\NormalTok{optional\_depends = moddrei}
\end{Highlighting}
\end{Shaded}

\section{Mod-Pakete}\label{mod-pakete}

Mods lassen sich zu Mod-Paketen\index{mod-paket@Mod-Paket} (Modpacks)
zusammenfassen, sodass mehrere Mods gemeinsam verpackt und verschoben
werden können. Das ist nützlich, wenn du einem Spieler mehrere Mods
bereitstellen möchtest, ohne dass er jeden einzeln herunterladen muss.

\begin{Shaded}
\begin{Highlighting}[]
\NormalTok{modpack1}
\NormalTok{├── modpack.conf (erforderlich, kennzeichnet Mod{-}Paket)}
\NormalTok{├── mod1}
\NormalTok{│   └── ... Mod{-}Dateien}
\NormalTok{└── mymod (optional)}
\NormalTok{    └── ... Mod{-}Dateien}
\end{Highlighting}
\end{Shaded}

Beachte, dass ein Mod-Paket kein \emph{Spiel} ist. Spiele haben eine
eigene Verzeichnisstruktur, die im Kapitel „Spiele`` erklärt wird.
